\documentclass{article}

\usepackage{hyperref}
\usepackage{graphicx}
\usepackage{fancyhdr}
\usepackage{extramarks}
\usepackage{amsmath}
\usepackage{cleveref}
\usepackage{amsthm}
\usepackage{amsfonts}
\usepackage{tikz}
\usepackage[plain]{algorithm}
\usepackage{algpseudocode}

\usetikzlibrary{automata,positioning}

%
% Basic Document Settings
%

\topmargin=-0.45in
\evensidemargin=0in
\oddsidemargin=0in
\textwidth=6.5in
\textheight=9.0in
\headsep=0.25in

\linespread{1.1}

\pagestyle{fancy}
\lhead{\hmwkAuthorName}
\chead{\hmwkClass , \hmwkClassInstructor: \hmwkTitle}
\rhead{\firstxmark}
\lfoot{\lastxmark}
\cfoot{\thepage}

\renewcommand\headrulewidth{0.4pt}
\renewcommand\footrulewidth{0.4pt}

\setlength\parindent{0pt}

%
% Create Problem Sections
%

\newcommand{\enterProblemHeader}[1]{
    \nobreak\extramarks{}{Problem \arabic{#1} continued on next page\ldots}\nobreak{}
    \nobreak\extramarks{Problem \arabic{#1} (continued)}{Problem \arabic{#1} continued on next page\ldots}\nobreak{}
}

\newcommand{\exitProblemHeader}[1]{
    \nobreak\extramarks{Problem \arabic{#1} (continued)}{Problem \arabic{#1} continued on next page\ldots}\nobreak{}
    \stepcounter{#1}
    \nobreak\extramarks{Problem \arabic{#1}}{}\nobreak{}
}

\setcounter{secnumdepth}{0}
\newcounter{partCounter}
\newcounter{homeworkProblemCounter}
\setcounter{homeworkProblemCounter}{1}
\nobreak\extramarks{Problem \arabic{homeworkProblemCounter}}{}\nobreak{}

%
% Homework Problem Environment
%
% This environment takes an optional argument. When given, it will adjust the
% problem counter. This is useful for when the problems given for your
% assignment aren't sequential. See the last 3 problems of this template for an
% example.
%
\newenvironment{homeworkProblem}[1][-1]{
    \ifnum#1>0
        \setcounter{homeworkProblemCounter}{#1}
    \fi
    \section{Problem \arabic{homeworkProblemCounter}}
    \setcounter{partCounter}{1}
    \enterProblemHeader{homeworkProblemCounter}
}{
    \exitProblemHeader{homeworkProblemCounter}
}

%
% Homework Details
%   - Title
%   - Due date
%   - Class
%   - Section/Time
%   - Instructor
%   - Author
%

\newcommand{\hmwkTitle}{Problem Set 2}
\newcommand{\hmwkClass}{CS229}
\newcommand{\hmwkClassInstructor}{Prof. Andrew Ng}
\newcommand{\hmwkAuthorName}{\textbf{Connor Yan}}

%
% Title Page
%

\title{
    \vspace{2in}
    \textmd{\textbf{\hmwkClass}}\\
    \textmd{\textbf{\hmwkTitle}}\\
    \vspace{3in}
}

\author{\hmwkAuthorName}
\date{\today}

\renewcommand{\part}[1]{\textbf{\large Part \Alph{partCounter}}\stepcounter{partCounter}\\}

% Alias for the Solution section header
\newcommand{\solution}{\textbf{\large Solution}}

% Probability commands: Expectation, Variance, Covariance, Bias
\newcommand{\E}{\mathbb{E}}
\newcommand{\Var}{\mathrm{Var}}
\newcommand{\Cov}{\mathrm{Cov}}
\newcommand{\Bias}{\mathrm{Bias}}
\newcommand{\dd}{\mathrm{d}}
\newcommand{\R}{\mathbb{R}}
\newcommand{\mb}{\mathrm}

\begin{document}

\maketitle

\pagebreak

\begin{homeworkProblem}
\begin{enumerate}
    \item[(a)] Regression on dataset B does not converge.
    
    \item[(b)] 
    \begin{figure}[!htbp]
        \centering
        \includegraphics[width=0.5\linewidth]{PS2/output/p01_ds1a.png}
        \caption{Dataset A}
        \label{fig:dsA}
    \end{figure}
    \begin{figure}[!htbp]
        \centering
        \includegraphics[width=0.5\linewidth]{PS2/output/p01_ds1b.png}
        \caption{Dataset B}
        \label{fig:dsB}
    \end{figure}

    Notice that dataset B is linearly separable, whereas dataset A is not. The loss function of this logistic regression is
    \[
        \ell(\theta)
        =-\frac{1}{m}\sum_{i=1}^m{\log \left( \frac{1}{1+e^{-y^{(i)}\theta^Tx^{(i)}}} \right) }
        =\frac{1}{m}\sum_{i=1}^m{\log \left( {1+e^{-y^{(i)}\theta^Tx^{(i)}}} \right) }
    \]
    As a negative log-likelihood, this is definitely greater than or equal to zero. However, for linear separable cases, if the fetched $\theta$ is good enough to rightly separate all data points, then this $\theta$ can be scaled up by an arbitrary positive number, and the classification is still right, with a smaller loss function. As $||\theta||\rightarrow\infty$,
    \[
        \ell(\theta)\rightarrow
        \frac{1}{m}\sum_{i=1}^m{\log \left( {1+e^{-\infty}} \right) }\rightarrow
        0,
    \]
    which means there is not a finite $\theta$ that yields to the minimum of $\ell(\theta)$.
    
    For non-linear-separable cases, where $y^{(i)}\theta^Tx^{(i)}<0$, terms in the loss function corresponding to data points that cannot be rightly classified are
    \[
        \log(1+e^{-y^{(i)}\theta^Tx^{(i)}}),
    \]
    As $||\theta||\rightarrow\infty$, 
    \[
        \log(1+e^{-y^{(i)}\theta^Tx^{(i)}})\rightarrow\infty
    \]
    This indicates there exists a $\theta$ that gives the minimum of the loss function, namely, the gradient descent can converge.

    \item[(c)] 

    \begin{enumerate}
        \item[i.] No. This is just a part of the gradient descent and does nothing with whether the minimum of the loss function exists.
    
        \item[ii.] No. This is just a part of the gradient descent and does nothing with whether the minimum of the loss function exists.
    
        \item[iii.] No. This change will not affect whether the dataset is linearly separable or not.
    
        \item[iv.] Yes. This can prevent infinitely scaling $\theta$ up to decrease the loss function.
    
        \item[v.] Probably would work if the modified dataset becomes non-linear-separable.

    \end{enumerate}

    \item[(d)] No. SVMs minimize $\frac{1}{2}||w||^2$, which punishes inifinitely scaling up of $w$.
    
\end{enumerate}
\end{homeworkProblem}

\begin{homeworkProblem}
    
\end{homeworkProblem}

\end{document}
